\documentclass[12pt]{article}

% Packages
\usepackage[utf8]{inputenc}
\usepackage[T1]{fontenc}
\usepackage{amsmath}
\usepackage{amssymb}
\usepackage{amsthm}
\input{common-defs}
\usepackage{graphicx}
\usepackage{natbib}
\input{typography-preamble}

% Title information
\title{Interpretable causal inference for high-stakes settings with optimal decision trees}
\author{Denis Agniel}
\date{\today}

\begin{document}

\maketitle

\begin{abstract}
% Placeholder: (1) Problem and gap---interpretable nuisances in DML without rate guarantees. (2) Approach---optimal trees in DML for ATT plus Rashomon intersection for a single tree per nuisance. (3) Main result---sufficient conditions for $o(n^{-1/4})$ and valid inference. (4) Empirical takeaway---competitive performance and interpretability; ATE discussed as extension in Discussion.
\end{abstract}

\section{Introduction}

Many consequential policy decisions rely on statistical models whose outputs directly affect financial payments, public reporting, and organizational reputation. Medicare quality measurement programs---including the Hospital Readmissions Reduction Program, Hospital Value-Based Purchasing, and Medicare Advantage Star Ratings---use risk-adjusted outcome models to evaluate provider performance and allocate billions of dollars in incentive payments each year \citep{cms_quality}. Because these models determine penalties, bonuses, and public rankings, they operate in a high-stakes regulatory environment where methodological choices must withstand scrutiny from providers, regulators, and the public. In such settings, interpretability is not merely a desirable feature but often a practical requirement for adoption: stakeholders demand models whose structure can be explained, audited, and defended, particularly when results may trigger financial penalties or allegations of poor care quality \citep{obermeyer2019, rajkomar2019, rudin2019}. Historically, this has led agencies such as the Centers for Medicare \& Medicaid Services to rely heavily on relatively transparent approaches, even when more flexible machine learning methods might offer improved predictive accuracy or better risk adjustment.

Estimation of causal effects in high-dimensional and flexible modeling environments has been greatly advanced by Double Machine Learning (DML), a semiparametric framework that combines machine learning with orthogonalized score functions to produce asymptotically normal, root-$n$ consistent estimates even when complex models are used to estimate nuisance functions \citep{chernozhukov2018}. A key technical requirement is that nuisance estimators converge sufficiently fast---typically faster than $n^{-1/4}$ in mean squared error---so that their estimation error is effectively second-order in the orthogonal identifying equations. In practice, off-the-shelf machine learning methods such as random forests, gradient boosting, and neural networks are commonly employed within DML pipelines \citep{bach2021, susmann2025}, including for healthcare provider evaluation. While these flexible learners often achieve good predictive performance, they are generally black-box models whose internal structure cannot be easily interpreted. Interpretability can be critical in many application domains where stakeholders demand transparent models that offer clear decision rules or partitions of the covariate space.

Optimal decision trees---including recent advances such as Generalized and Scalable Optimal Sparse Decision Trees (GOSDT)---offer a compelling trade-off: they are fully interpretable, provide globally optimized decision rules, and avoid some of the pathologies of greedy tree induction \citep{lin2020}. These optimal trees have been studied from a computational and predictive performance perspective, but their statistical properties in semiparametric inference contexts remain largely unexplored. Despite the intuitive appeal of combining interpretable tree models with DML, there is a noticeable absence of formal results in the literature establishing that optimal decision tree estimators satisfy the rate conditions (e.g., $o(n^{-1/4})$ convergence) required for valid inference in the DML framework. Most existing work using trees in DML settings relies on empirical evidence or practitioner heuristics rather than rigorous guarantees; moreover, statistical convergence analyses for tree-based estimators traditionally focus on greedy algorithms or ensembles, and only under specific smoothness and complexity constraints.

In this paper, we address this gap. We propose incorporating optimal decision trees as the nuisance function learners in a Double Machine Learning procedure and investigate the theoretical conditions under which these trees attain the necessary convergence rates to preserve root-$n$ consistency and asymptotic normality of the target causal parameter estimate. Our contributions are as follows. We introduce a DML framework where optimal, interpretable decision trees (e.g., GOSDT) are used to estimate nuisance functions in causal models. We derive sufficient conditions on tree complexity and model regularization under which these tree estimators satisfy mean squared error rates faster than $n^{-1/4}$ in relevant function classes. We provide empirical evidence, including simulation studies and real data examples, demonstrating that interpretable trees within DML achieve competitive estimation quality while offering transparent decision rules. By bridging the gap between interpretable machine learning and semiparametric causal inference, this work aims to develop methods that are transparent enough for regulatory acceptance, flexible enough to capture complex risk structure, and theoretically grounded enough to support valid inference---particularly important for modern quality measurement, where methodological credibility is inseparable from stakeholder trust.

The rest of the paper is organized as follows. Section~\ref{sec:methods} presents the DML framework with optimal tree nuisance learners and states the main theoretical results. Section~\ref{sec:rashomon} describes the use of Rashomon sets to obtain a single tree per nuisance across cross-fitting folds. Section~\ref{sec:results} reports simulation and empirical findings. Section~\ref{sec:discussion} discusses limitations and extensions, and Section~\ref{sec:conclusion} concludes.

\section{Methods}
\label{sec:methods}

\subsection{Setting and target parameter}

We observe \iid data $O_1, \ldots, O_n$, where each $O = (X, A, Y)$ consists of a covariate vector $X \in \cX$, a binary treatment $A \in \{0,1\}$, and an outcome $Y$. Let $Y(0)$ and $Y(1)$ denote the potential outcomes under control and treatment, so that $Y = A Y(1) + (1-A) Y(0)$. Our target parameter is the average treatment effect on the treated (ATT),
\begin{equation}
\label{eq:att}
\theta_0 = \E[Y(1) - Y(0) \mid A=1].
\end{equation}

We assume unconfoundedness (selection on observables): $(Y(0), Y(1)) \perp A \mid X$. We also assume overlap: $0 < \Prob(A=1 \mid X) < 1$ almost surely (or at least on the support of $X$ among the treated). Under these conditions, the ATT is identified via the propensity score $e(X) = \Prob(A=1 \mid X)$ and the outcome regressions $m(a,X) = \E[Y \mid A=a, X]$ for $a \in \{0,1\}$. The nuisance functions required for estimation are $e(X)$, $m(0,X)$, and $m(1,X)$; we collect them in the vector $\eta = (e, m_0, m_1)$ with $m_0(X) = m(0,X)$ and $m_1(X) = m(1,X)$.

The ATT can be estimated using a Neyman-orthogonal (efficient influence) score $\psi(O; \theta, \eta)$ that satisfies $\E[\psi(O; \theta_0, \eta_0)] = 0$ and has zero first-order Gateaux derivative with respect to $\eta$ at $\eta_0$, so that first-order nuisance estimation error does not affect the limiting distribution of the estimator \citep{chernozhukov2018}. Let $\pi = \Prob(A=1)$. A standard form for the orthogonal score for the ATT is
\begin{equation}
\label{eq:score}
\psi(O; \theta, \eta) = \frac{A}{\pi}\bigl( Y - m_0(X) - \theta \bigr) - \frac{1}{\pi} \, \frac{A - e(X)}{1 - e(X)} \bigl( m_1(X) - m_0(X) \bigr).
\end{equation}
At the true $\theta_0$ and $\eta_0$, $\E[\psi(O; \theta_0, \eta_0)] = 0$ (the second term has mean zero by the propensity score, and the first has mean $\theta_0 - \theta$); the estimator $\hat\theta$ will be defined by solving the empirical (cross-fitted) version of this moment condition.

\subsection{Double Machine Learning with orthogonal scores}

Double Machine Learning (DML) estimates $\theta_0$ by combining sample splitting with flexible nuisance estimation so that the orthogonal score can be evaluated with out-of-sample nuisance predictions, avoiding overfitting bias \citep{chernozhukov2018}. The sample is split into $K$ folds. For each fold $k$, nuisance functions $\hat\eta^{(-k)} = (\hat{e}^{(-k)}, \hat{m}_0^{(-k)}, \hat{m}_1^{(-k)})$ are estimated using only data not in fold $k$. On fold $k$, the score $\psi(O; \theta, \hat\eta^{(-k)})$ is evaluated. The DML estimator $\hat\theta$ solves the aggregated moment condition
\begin{equation}
\label{eq:dml}
\frac{1}{n} \sum_{i=1}^n \psi\bigl(O_i; \hat\theta, \hat\eta^{(-k(i))}\bigr) = 0,
\end{equation}
where $k(i)$ denotes the fold containing observation $i$.

For $\hat\theta$ to be $\sqrt{n}$-consistent and asymptotically normal, the nuisance estimators must converge sufficiently fast. Specifically, \citet{chernozhukov2018} show that if the nuisance estimators achieve mean squared error ($\MSE$) rate $o(n^{-1/4})$ in the relevant norms, then the linearization remainder in the expansion of $\hat\theta$ is $o_p(n^{-1/2})$, so that $\sqrt{n}(\hat\theta - \theta_0)$ is asymptotically normal with variance equal to that of the influence function. The $n^{-1/4}$ rate arises because the remainder typically involves products of nuisance errors; to ensure these products are $o_p(n^{-1/2})$, each nuisance error must be $o_p(n^{-1/4})$. More generally, \emph{any} fold-specific nuisance estimators that are $o_p(n^{-1/4})$-close to the true nuisances (in the norms appearing in the score expansion) yield the same conclusion; we state this as a general theorem in the Appendix (Section~\ref{app:alt-nuisances}). In the next subsection we specialize to optimal decision trees for estimating $\eta$.

\subsection{Optimal decision trees as nuisance learners}

We estimate the nuisance functions $e$, $m_0$, and $m_1$ using optimal decision trees. A decision tree partitions the covariate space $\cX$ into regions (leaves) and assigns a constant prediction to each region. Optimal trees, such as those produced by the Generalized and Scalable Optimal Sparse Decision Trees (GOSDT) algorithm \citep{lin2020}, minimize empirical risk plus a complexity penalty (e.g., on the number of leaves or tree depth), yielding globally optimized partitions and predictions. Trees are flexible in both the partition structure and the leaf values, and they yield interpretable rules.

We use a single \emph{loss} $L$ (e.g., squared error for continuous outcomes, cross-entropy for binary outcomes) and a complexity penalty for all nuisances. For the propensity score $e(X)$, we fit a tree that predicts $A$ given $X$; for the outcome regressions $m_0(X)$ and $m_1(X)$, we fit separate trees on units with $A=0$ and $A=1$ to predict $Y$ given $X$. Thus we obtain separate trees for $\hat{e}$, $\hat{m}_0$, and $\hat{m}_1$.

For each cross-fitting fold $k$, we train $\hat{e}^{(-k)}$, $\hat{m}_0^{(-k)}$, and $\hat{m}_1^{(-k)}$ on the training sample (the complement of fold $k$) using the chosen tree learner. The predictions $\hat\eta^{(-k)}(X) = (\hat{e}^{(-k)}(X), \hat{m}_0^{(-k)}(X), \hat{m}_1^{(-k)}(X))$ are then used in the score $\psi(O; \theta, \hat\eta^{(-k)})$ for observations in fold $k$. The fitted $\hat\eta$ thus consists of optimal-tree nuisance learners that are interpretable and that we will show can satisfy the rate conditions required for valid inference.

\subsection{Main result}

We treat the nuisance estimators as \emph{penalized empirical risk minimizers}: they minimize empirical risk plus a penalty $\lambda_n \times (\text{number of leaves})$ over axis-aligned decision trees. The penalty controls complexity; the theoretical rate index is $s_n \asymp n^{d/(2\beta+d)}$. In practice we use a GOSDT-style solver that minimizes this objective; the same rates apply if the solver's suboptimality is $o(s_n \log n / n)$ (see Appendix).

Under the following conditions, these tree nuisance estimators achieve the $o(n^{-1/4})$ $\MSE$ rate required for DML, so that the resulting DML estimator of the ATT is $\sqrt{n}$-consistent and asymptotically normal.

\begin{itemize}
\item \textbf{Identification:} Unconfoundedness and overlap hold; outcomes and the propensity score are bounded so that the relevant moments exist.
\item \textbf{Approximation:} The support of $X$ is contained in $[0,1]^d$ (or a compact set, after transformation). The true nuisances $\eta_0 = (e_0, m_{0,0}, m_{1,0})$ lie in a H\"older class $H^\beta([0,1]^d)$ for some $\beta > 0$.
\item \textbf{Loss:} The loss $L$ is such that excess risk $R(\hat\eta) - R(\eta_0)$ controls the $L^2(P)$ norm in the sense required by the DML expansion (e.g., squared error, or cross-entropy when $Y$ is binary).
\item \textbf{Complexity:} The tree complexity (number of leaves) satisfies $s_n \asymp n^{d/(2\beta+d)}$.
\item \textbf{DML rate:} The smoothness and dimension satisfy $\beta/(2\beta+d) > 1/4$, or equivalently $\beta > d/2$, so that the tree rate is faster than $n^{-1/4}$.
\end{itemize}
Under these conditions, the penalized-ERM tree nuisance estimators $\hat{e}$, $\hat{m}_0$, and $\hat{m}_1$ (trained per fold as above) achieve $\|\hat\eta - \eta_0\|_{L^2(P)} = O_p(n^{-\beta/(2\beta+d)})$, hence $o_p(n^{-1/4})$, by the excess-risk rate and the loss--norm link (see Appendix). It follows from standard DML theory \citep{chernozhukov2018} that the estimator $\hat\theta$ defined by \eqref{eq:dml} satisfies $\sqrt{n}(\hat\theta - \theta_0) \leadsto N(0, \sigma^2)$, where $\sigma^2$ is the variance of the influence function $\psi(O; \theta_0, \eta_0)$. Formal statements and proof details are given in the Appendix.

In high dimensions, the complexity of the tree class can grow quickly; to retain the same rate, one may impose sparsity of the regression function, restricted splitting depth, or margin conditions (see Appendix).

\subsection{Estimation and inference}

In practice we use $K$-fold cross-fitting with $K = 5$ or $K = 10$ (or another user-specified value). The regularization parameter $\lambda$ (penalty per leaf) can be chosen by cross-validation: for each nuisance, minimize held-out loss (the chosen loss $L$ on a validation split) over a grid of $\lambda$ values. A theory-driven default is $\lambda \propto (\log n)/n$ (e.g.\ $\lambda = (\log n)/n$), which matches the rate conditions in the Appendix and typically yields trees with $O(n^{d/(2\beta+d)})$ leaves under the conditions of the theorem; the proportionality constant can be tuned by cross-validation if desired. Sensitivity to the choice of $\lambda$ can be examined in simulations.

The asymptotic variance $\sigma^2$ of $\sqrt{n}(\hat\theta - \theta_0)$ is the variance of $\psi(O; \theta_0, \eta_0)$. We estimate it by the empirical variance of the estimated influence function,
\begin{equation}
\hat\sigma^2 = \frac{1}{n} \sum_{i=1}^n \psi\bigl(O_i; \hat\theta, \hat\eta^{(-k(i))}\bigr)^2.
\end{equation}
A Wald-type 95\% confidence interval for $\theta_0$ is
\begin{equation}
\hat\theta \pm 1.96 \, \frac{\hat\sigma}{\sqrt{n}}.
\end{equation}

\section{Rashomon sets: single tree per nuisance}
\label{sec:rashomon}

\subsection{Rashomon set per fold}

For each cross-fitting fold $k$, we fit optimal trees on the training part (data not in fold $k$) as in Section~\ref{sec:methods}. The \emph{Rashomon set} for fold $k$ is the set of trees that are ``almost as good'' as the best tree on that training set: trees (in $\cT_{M_n}$ or $\cT_{s_n}$) whose empirical penalized risk on $\mathcal{D}^{(-k)}$ is within a factor $(1+\varepsilon)$ of the minimum, or within an additive $\varepsilon$ of the best. Formally, $R_k(\varepsilon) = \{ f \in \cT_{M_n} : R_n^{(-k)}(f) + \lambda_n \cdot (\#\text{leaves of } f) \le (1+\varepsilon) \cdot \text{best}^{(-k)} \}$, where $\text{best}^{(-k)}$ is the minimal penalized risk on the training data for fold $k$ (see Appendix~\ref{app:rashomon}). Thus we obtain multiple trees that are good enough on that fold, rather than a single optimal tree. Algorithms such as SPLIT \citep{babbarNearOptimalDecisionTrees2025} or GOSDT-style solvers can enumerate or approximate the Rashomon set for decision trees; methods that enumerate the full Rashomon set for sparse trees have been studied by \citet{xinExploringWholeRashomon2022}.

\subsection{Intersection across folds}

The goal is to select a \emph{single} tree per nuisance (one propensity tree and one or two outcome trees for $m_0$, $m_1$) for interpretability and reporting, while preserving valid inference. To that end, we define the \emph{intersection} of the Rashomon sets across folds: the set of tree \emph{structures} (splits) that appear in every fold's Rashomon set when each fold's trees are fitted on their respective training data $\mathcal{D}^{(-k)}$. A structure in the intersection is one that, when refit on each fold's training data, has penalized risk within $(1+\varepsilon)$ of the best for that fold. One can compute the intersection by enumerating or approximating the Rashomon sets per fold and matching structures (e.g., via a canonical representation or JSON). We then choose one structure from the intersection per nuisance (e.g., the smallest or the first).

For use in the DML estimator, for each fold $k$ we take $\tilde{\eta}^{(-k)}$ to be the tree from the intersection \emph{refit on $\mathcal{D}^{(-k)}$}---or equivalently, the fold-$k$ instance of that structure from the Rashomon set of fold $k$. The score $\psi(O_i; \theta, \tilde{\eta}^{(-k(i))})$ is evaluated with these fold-specific fitted nuisances, so that the orthogonality and cross-fitting structure of DML are preserved. The output is one reportable propensity tree and one reportable outcome tree (or one structure each for $m_0$ and $m_1$) for stakeholders, with fold-specific fitted versions used in the estimator.

\subsection{Rationale and caveats}

In high-stakes settings, a single stable tree per nuisance is more interpretable and auditable than $K$ fold-specific trees that may differ across folds. Stakeholders can inspect one propensity tree and one outcome tree rather than $K$ versions of each.

From a theoretical standpoint, under the conditions of the Appendix (Definition~\ref{def:rashomon-fold}, Definition~\ref{def:intersection}, Proposition~\ref{prop:rashomon-dml}), choosing nuisances from the Rashomon set---or from the intersection, with refit per fold---yields $\tilde{\eta}^{(-k)}$ that are $o_p(n^{-1/4})$-close to $\eta_0$. The general theorem (Appendix~\ref{app:alt-nuisances}, Theorem~\ref{thm:dml-alt}) then implies that the DML estimator with these nuisances remains $\sqrt{n}$-consistent and asymptotically normal with the same limiting variance.

A caveat is that the intersection tree may have higher bias or a larger constant than a fold-specific optimal tree, because it must perform well on every fold's training set. This bias--interpretability trade-off can be examined in simulations (Section~\ref{sec:results}), where we compare DML with fold-specific optimal trees versus DML with the Rashomon intersection.

\section{Results}
\label{sec:results}

\subsection{Simulation design}
% DGPs: general causal model (not necessarily partially linear), nonlinear nuisances; at least one smooth, one piecewise/heterogeneous; design identifies ATT.
% Sample sizes: e.g., $n \in \{500, 1000, 2000, 5000\}$.
% Comparisons: (i) DML with optimal trees per fold, (ii) DML with Rashomon intersection, (iii) DML with random forest/gradient boosting, (iv) DML with linear/logistic nuisances.
% Target: ATT. Metrics: bias, RMSE, coverage, 95\% CI length; optionally nuisance MSE. Replications: e.g., 500 or 1000.

\subsection{Simulation results}
% Tables/figures: bias, RMSE, coverage by method and $n$ for ATT.
% Narrative: tree-DML comparable (or better in heterogeneous DGPs) to forest-DML; linear DML suffers when nuisances nonlinear.
% Rashomon intersection: compare fold-specific vs single-tree estimator; document bias/efficiency trade-off; discuss acceptability of bias for interpretability gain.
% If available: nuisance MSE decay with $n$ (consistent with $o(n^{-1/4})$).

\subsection{Real-data example: quality measurement for SUD treatment facilities}
% Application: quality measurement for substance use disorder (SUD) treatment facilities (risk-adjusted outcomes or treatment effects, facility- or patient-level).
% Nuisances: optimal trees; optionally Rashomon intersection for single reportable tree.
% Report: point estimate and CI for ATT; brief comparison to black-box DML if feasible.
% Interpretability: show one or two fitted trees (or rules); explain partition of covariate space (e.g., facility/patient characteristics defining risk subgroups).

\subsection{Interpretability illustration}
% If not folded into previous subsection: one figure or table summarizing transparent decision rules (tree diagram or list of rules) and stakeholder communication in SUD facility context.

\section{Discussion}
\label{sec:discussion}

\subsection{Limitations}
% Function class/assumptions: nuisances well approximated by trees of bounded complexity; impact of misspecification.
% Computation: cost of optimal tree fitting vs greedy/forests; scalability.
% Choice of complexity: cross-validation/tuning; sensitivity of inference.
% Rashomon intersection: lack of formal theory; possible bias vs interpretability trade-off.

\subsection{Extensions}
% ATE: same DML-with-trees and Rashomon approach applies; orthogonal score and nuisances differ; mention briefly.
% Other parameters (CATE, LATE) or identifying assumptions; other interpretable learners; high-dimensional $X$; formal theory for Rashomon-intersection if feasible.

\subsection{Implications for high-stakes settings}
A central contribution of this work is to show that \emph{interpretable} learners (decision trees) can satisfy the DML nuisance-rate requirement and thus deliver valid inference---transparent models need not sacrifice statistical validity. Tree-DML is a candidate for settings like Medicare quality measurement: interpretable nuisances, valid confidence intervals, and transparent rules. Caveats include the need for validation, domain input, and regulatory alignment.

\section{Conclusion}
\label{sec:conclusion}
% Paragraph 1: Restate gap (interpretable nuisances in DML without rate guarantees) and what we did: optimal trees as nuisances in DML for ATT, sufficient conditions for $o(n^{-1/4})$, Rashomon intersection for single tree per nuisance, empirical evidence of competitive performance and interpretability.
% Paragraph 2: Broader impact---transparent, statistically principled causal estimation for high-stakes applications (policy, healthcare) without sacrificing validity.

\appendix
\section{Proof of the nuisance rate}
\label{app:nuisance-rate}

We give formal statements and proof ideas for the tree nuisance rate. The setting is that of a generic nuisance (e.g., regression of $Y$ on $X$; propensity is analogous with binary outcome). Cross-fitting is handled as in the main text (train on the complement of each fold). The tree class $\cT_s$ is the set of axis-aligned decision trees with at most $s$ leaves (functions constant on each leaf). We take $\cT_{s_n}$ as a sieve indexed by $n$. The covariate space is $[0,1]^d$ (or a compact set after transformation). The norm $\|\cdot\|_{L^2(P)}$ denotes $L^2(P)$ with respect to the distribution of $X$.

Let $L(y, f(x))$ be a loss function. Define the risk $R(f) = \E[L(Y, f(X))]$ and the excess risk $R(f) - R(\eta_0)$. The estimator is the penalized empirical risk minimizer:
\begin{equation}
\label{eq:perm}
\hat\eta \in \argmin_{f \in \cT_{M_n}} \Bigl[ \frac{1}{n}\sum_{i=1}^n L(Y_i, f(X_i)) + \lambda_n \cdot (\#\text{leaves of } f) \Bigr],
\end{equation}
where $\cT_{M_n}$ is the set of axis-aligned decision trees with at most $M_n$ leaves, with $M_n \to \infty$ and $M_n \log n / n \to 0$ as $n \to \infty$ (e.g.\ $M_n = n$ or $M_n \propto s_n \log n$). This matches the objective of a GOSDT-style solver (loss plus regularization times number of leaves). The constant $s_n$ below is the theoretical rate index, $s_n \asymp n^{d/(2\beta+d)}$. In practice, $\lambda_n$ may be taken proportional to $(\log n)/n$ (e.g.\ $\lambda_n = c (\log n)/n$ for a constant $c$) to match the rate conditions; the constant $c$ can be chosen by cross-validation. We assume a \textbf{loss--norm link}: there exists a function $\phi(\delta) \to 0$ as $\delta \to 0$ such that $R(f) - R(\eta_0) \le \delta$ implies $\|f - \eta_0\|_{L^2(P)}^2 \le \phi(\delta)$. This holds for squared error ($\phi(\delta)=\delta$) and for cross-entropy with bounded probabilities (standard bounds). Any procedure that attains this minimum or an $\varepsilon_n$-optimal solution (penalized risk $\le \inf_f (\cdots) + \varepsilon_n$) fits the framework provided $\varepsilon_n = o(s_n \log n / n)$; GOSDT-style solvers minimize this objective (optionally with a time or tolerance limit).

\begin{lemma}[Approximation by tree partitions]
\label{lem:approx}
Let $\cT_s$ be the class of axis-aligned trees with at most $s$ leaves. If $\eta_0 \in H^\beta([0,1]^d)$, then
\[
\inf_{f \in \cT_s} \bigl[ R(f) - R(\eta_0) \bigr] \le C s^{-2\beta/d}.
\]
\end{lemma}
\begin{proof}
We first give precise notation, then prove the bound for squared-error loss, and extend to generic loss.

\paragraph{Notation.}
For $\beta > 0$, the H\"older class $H^\beta([0,1]^d)$ is defined as follows. Let $\ell = \lfloor \beta \rfloor$ and $\gamma = \beta - \ell \in (0,1]$. A (bounded) function $\eta$ on $[0,1]^d$ belongs to $H^\beta([0,1]^d)$ if (i) all partial derivatives of $\eta$ of order $k$ exist and are bounded for $k = 0, \ldots, \ell$, and (ii) for every multi-index $\alpha$ with $|\alpha| = \ell$, the derivative $D^\alpha \eta$ is $\gamma$-H\"older continuous: $|D^\alpha \eta(x) - D^\alpha \eta(x')| \le C |x - x'|^\gamma$ for all $x, x' \in [0,1]^d$. The norm $\|\eta\|_{H^\beta}$ can be taken as the maximum over the relevant derivative bounds and H\"older constants. For $\beta \in (0,1]$, this reduces to the usual condition $|\eta(x) - \eta(x')| \le \|\eta\|_{H^\beta} |x - x'|^\beta$. Every axis-aligned partition of $[0,1]^d$ into at most $s$ cells can be realized by some $f \in \cT_s$ (piecewise-constant on the partition).

\paragraph{Step 1: Partition.}
Let $k = \lfloor s^{1/d} \rfloor \ge 1$. Divide each coordinate into $k$ intervals of length $1/k$, yielding $N = k^d \le s$ axis-aligned cells $A_1, \ldots, A_N$ that partition $[0,1]^d$. Each cell has diameter $\diam(A_j) \le \sqrt{d}/k \le C_d s^{-1/d}$.

\paragraph{Step 2: Approximant.}
Define $f_s$ by: on each cell $A_j$, set $f_s(x) = \E[\eta_0(X) \mid X \in A_j]$ (or $\int_{A_j} \eta_0 \, dP / P(A_j)$ when $P(A_j) > 0$). Then $f_s$ is piecewise constant on this partition and belongs to $\cT_s$.

\paragraph{Step 3: Bound for squared error.}
For $L(y,f) = (y-f)^2$, excess risk equals $\|f - \eta_0\|_{L^2(P)}^2$. We have
\[
\|f_s - \eta_0\|_{L^2(P)}^2 = \sum_{j=1}^N \int_{A_j} (f_s - \eta_0)^2 \, dP.
\]
On $A_j$, $f_s$ is the $L^2(P)$-projection of $\eta_0$ onto constants, so $\int_{A_j} (f_s - \eta_0)^2 \, dP \le \int_{A_j} (\eta_0 - c)^2 \, dP$ for any constant $c$. Choose $c = \eta_0(x_j)$ for some $x_j \in A_j$. By the H\"older condition, $|\eta_0(x) - \eta_0(x_j)| \le \|\eta_0\|_{H^\beta} |x - x_j|^\beta \le \|\eta_0\|_{H^\beta} (\diam(A_j))^\beta$. Hence
\[
\int_{A_j} (\eta_0 - \eta_0(x_j))^2 \, dP \le P(A_j) \cdot (\diam(A_j))^{2\beta} \|\eta_0\|_{H^\beta}^2 \le P(A_j) \cdot C s^{-2\beta/d},
\]
where $C$ absorbs $\|\eta_0\|_{H^\beta}^2$, $d$, and the diameter bound. Summing over $j$ and using $\sum_j P(A_j) = 1$ gives $\|f_s - \eta_0\|_{L^2(P)}^2 \le C s^{-2\beta/d}$. Thus $\inf_{f \in \cT_s} \|f - \eta_0\|_{L^2(P)}^2 \le \|f_s - \eta_0\|_{L^2(P)}^2 \le C s^{-2\beta/d}$.

\paragraph{Step 4: Squared error conclusion.}
For squared error, $R(f_s) - R(\eta_0) = \|f_s - \eta_0\|_{L^2(P)}^2 \le C s^{-2\beta/d}$, so the lemma holds.

\paragraph{Generic loss.}
Under the loss--norm link, any $f_s$ with $\|f_s - \eta_0\|_{L^2(P)}^2 \le C s^{-2\beta/d}$ satisfies $R(f_s) - R(\eta_0) \le \phi(C s^{-2\beta/d}) = O(s^{-2\beta/d})$ when $\phi(\delta) \lesssim \delta$ (e.g.\ squared error and cross-entropy with bounded probabilities). Thus $\inf_{f \in \cT_s} [R(f) - R(\eta_0)] \le C s^{-2\beta/d}$ for such losses. For further discussion of piecewise-constant approximation see \citet{gyorfiLeastSquaresEstimates2002} or \citet{devoreConstructiveApproximation1993}.

\paragraph{Case $\beta > 1$.}
For $\eta_0 \in H^\beta([0,1]^d)$ with $\beta > 1$, piecewise-constant approximation on a regular partition of $[0,1]^d$ into $O(s)$ cells of diameter $O(s^{-1/d})$ achieves $L^2$ error of order $s^{-\beta/d}$ (hence squared error $s^{-2\beta/d}$). This follows from the fact that on each cell the function can be approximated by its Taylor polynomial of degree $\ell = \lfloor \beta \rfloor$ at a point in the cell, with remainder bounded in $L^2$ by $O((\diam)^\beta)$; averaging to a constant on the cell yields the same rate. See, e.g., \citet{devoreConstructiveApproximation1993} (approximation by piecewise polynomials and piecewise constants) or \citet{gyorfiLeastSquaresEstimates2002} (rates for $H^\beta$ in nonparametric regression). Thus $\inf_{f \in \cT_s} \|f - \eta_0\|_{L^2(P)}^2 \le C s^{-2\beta/d}$, and by the loss--norm link, $\inf_{f \in \cT_s} [R(f) - R(\eta_0)] \le C s^{-2\beta/d}$.
\end{proof}

\begin{lemma}[Complexity of tree class]
\label{lem:complexity}
The covering number satisfies $\log N(\epsilon, \cT_s, L^2(P)) \lesssim s \log(1/\epsilon)$, or equivalently the VC dimension of $\cT_s$ is of order $s \log s$.
\end{lemma}
\begin{proof}
We define the relevant notions, bound the VC (pseudo-) dimension of $\cT_s$, and deduce the covering number bound.

\paragraph{Definitions.}
Assume functions in $\cT_s$ are uniformly bounded (e.g.\ in $[0,1]$), which holds for propensity and bounded outcomes. The \emph{covering number} $N(\epsilon, \cT_s, L^2(P))$ is the minimal number of closed $L^2(P)$-balls of radius $\epsilon$ whose union contains $\cT_s$. The \emph{VC dimension} of a class of sets is the largest $n$ such that some set of $n$ points is shattered. For real-valued functions, the \emph{pseudo-dimension} (VC subgraph dimension) is the VC dimension of the class of subgraphs. A standard fact (e.g.\ van der Vaart and Wellner, 1996, Theorem 2.6.7) is that for a uniformly bounded function class with VC subgraph dimension at most $d$, $\log N(\epsilon, \cT_s, L^2(P)) \le C d \log(1/\epsilon)$ for $\epsilon$ small enough.

\paragraph{VC bound.}
Each $f \in \cT_s$ is piecewise constant on an axis-aligned partition with at most $s$ cells; that is, $f = \sum_{j=1}^s a_j \mathbf{1}_{A_j}$ for disjoint axis-aligned cells $A_j$ and constants $a_j \in [0,1]$. The class of indicator functions of such cells (or of sets realizable as unions of at most $s$ axis-aligned boxes) has VC dimension $O(s \log s)$; see \citet{bartlettRademacherGaussianComplexities2002} or \citet{bartlettLocalRademacherComplexities2006} for decision trees and axis-aligned partitions. The class $\cT_s$ is contained in the set of linear combinations of $s$ such indicators with coefficients in $[0,1]$, so its pseudo-dimension is $O(s \log s)$.

\paragraph{Covering number.}
By the standard entropy bound cited above, $\log N(\epsilon, \cT_s, L^2(P)) \lesssim s \log s \cdot \log(1/\epsilon)$. Thus the VC (pseudo-) dimension of $\cT_s$ is of order $s \log s$, and the covering number bound above suffices for the oracle inequality in Lemma~\ref{lem:oracle} (the complexity term $s_n \log n / n$ is consistent with $\log N \lesssim s \log s \cdot \log(1/\epsilon)$ when $\epsilon \sim 1/\sqrt{n}$, giving $\log N \lesssim s \log n$ for $s \ge 2$). In Lemma~\ref{lem:oracle} we apply this with $s = O(s_n)$ (since $\hat\eta \in \cT_{O_p(s_n)}$ and $f_n \in \cT_{s_n}$ by Lemma~\ref{lem:hats}).
\end{proof}

\begin{lemma}[Number of leaves of the penalized minimizer]
\label{lem:hats}
Let $\hat{s} := \#\text{leaves of }\hat\eta$. Under $\lambda_n \asymp (\log n)/n$ and $s_n \asymp n^{d/(2\beta+d)}$, the penalized minimizer $\hat\eta$ in \eqref{eq:perm} satisfies $\hat{s} = O_p(s_n)$.
\end{lemma}
\begin{proof}[Proof idea]
Take $f_n \in \cT_{s_n}$ with $R(f_n) - R(\eta_0) \le \inf_{f \in \cT_{s_n}}[R(f)-R(\eta_0)] + o(1)$. By definition of $\hat\eta$, $R_n(\hat\eta) + \lambda_n \hat{s} \le R_n(f_n) + \lambda_n s_n$, so $\lambda_n \hat{s} \le R_n(f_n) - R_n(\hat\eta) + \lambda_n s_n \le R_n(f_n) + \lambda_n s_n$. Boundedness of $R(\cdot)$ and the usual empirical-process bound give $R_n(f_n) \le R(f_n) + O_p(\sqrt{s_n \log n / n}) = O(1) + O_p(\sqrt{s_n \log n / n})$. Thus $\hat{s} \le s_n + (R_n(f_n) - R_n(\hat\eta))/\lambda_n \le s_n + O_p(1/\lambda_n) + O_p(\sqrt{s_n \log n / n}/\lambda_n)$. With $\lambda_n \asymp (\log n)/n$, the second term is $O_p(n/\log n)$ and the third is $O_p(n \sqrt{s_n \log n / n} / \log n) = O_p(\sqrt{n s_n / \log n})$. For $s_n \asymp n^{d/(2\beta+d)}$, $\sqrt{n s_n}/\sqrt{\log n}$ is $o(s_n)$ when $d/(2\beta+d) < 1$ (which holds for $\beta > 0$), so $\hat{s} \le s_n + o_p(s_n) = O_p(s_n)$.
\end{proof}

\begin{lemma}[Oracle inequality for penalized ERM]
\label{lem:oracle}
For $\hat\eta$ the penalized ERM~\eqref{eq:perm} (minimize $R_n(f) + \lambda_n \cdot (\#\text{leaves of } f)$ over $\cT_{M_n}$) with $\lambda_n \asymp (\log n)/n$ and $s_n$ as above,
\[
R(\hat\eta) - R(\eta_0) \le C \Bigl( \inf_{f \in \cT_{s_n}} \bigl[ R(f) - R(\eta_0) \bigr] + \frac{s_n \log n}{n} \Bigr) + o_p(1).
\]
\end{lemma}
\begin{proof}
We assume the following in addition to the loss--norm link. \textbf{Boundedness:} $Y$ and all $f \in \cT_{M_n}$ take values in a bounded interval (e.g., $[0,1]$ for propensity or bounded outcomes); this is already implied by Lemma~\ref{lem:complexity}. \textbf{Loss:} $L(y,\cdot)$ is Lipschitz in the second argument uniformly in $y$, or quadratic (e.g., $|L(y,a)-L(y,b)| \le M|a-b|$ or $\le M(a-b)^2$); squared error and cross-entropy with bounded probabilities satisfy this. \textbf{Sieve:} $s_n \to \infty$ and $s_n \log n / n \to 0$ as $n \to \infty$. By Lemma~\ref{lem:hats}, $\hat{s} := \#\text{leaves of }\hat\eta = O_p(s_n)$, so $\hat\eta \in \cT_{C s_n}$ with high probability for some constant $C$.

\paragraph{Step 1: Decomposition.}
Let $R_n(f) = \frac{1}{n}\sum_{i=1}^n L(Y_i, f(X_i))$. By definition of $\hat\eta$, for $f_n \in \cT_{s_n}$ we have $R_n(\hat\eta) + \lambda_n \hat{s} \le R_n(f_n) + \lambda_n s_n$, so $R_n(\hat\eta) - R_n(f_n) \le \lambda_n(s_n - \hat{s}) \le \lambda_n s_n = O(s_n \log n / n)$. Fix $f_n \in \cT_{s_n}$ such that $R(f_n) - R(\eta_0) \le \inf_{f \in \cT_{s_n}}[R(f)-R(\eta_0)] + o(1)$. Then
\[
R(\hat\eta) - R(\eta_0) = \bigl(R(\hat\eta) - R_n(\hat\eta)\bigr) + \bigl(R_n(\hat\eta) - R_n(f_n)\bigr) + \bigl(R_n(f_n) - R(f_n)\bigr) + \bigl(R(f_n) - R(\eta_0)\bigr).
\]
The second bracket is $\le \lambda_n s_n = O(s_n \log n / n)$. Hence
\[
R(\hat\eta) - R(\eta_0) \le \underbrace{\bigl(R(\hat\eta) - R_n(\hat\eta)\bigr)}_{\text{(I)}} + \underbrace{\bigl(R_n(f_n) - R(f_n)\bigr)}_{\text{(II)}} + \underbrace{\bigl(R(f_n) - R(\eta_0)\bigr)}_{\text{(III)}} + O(s_n \log n / n).
\]

\paragraph{Step 2: Approximation term.}
Term (III) satisfies $R(f_n) - R(\eta_0) \le \inf_{f \in \cT_{s_n}}[R(f)-R(\eta_0)] + o(1)$. Thus (III) is at most the approximation term in the stated bound (up to a constant $C \ge 1$ and $o(1)$).

\paragraph{Step 3: Empirical process.}
Terms (I) and (II) are each of the form $(P_n - P)(L(Y,f(X)) - L(Y,\eta_0(X)))$ for some $f$ in a class of trees with $O_p(s_n)$ leaves ($\hat\eta$ and $f_n$). By Lemma~\ref{lem:complexity}, $\log N(\epsilon, \cT_s, L^2(P)) \lesssim s \log s \cdot \log(1/\epsilon)$. With $\hat{s} = O_p(s_n)$, the same entropy and maximal inequality give (I) and (II) together $O_p(s_n \log n / n)$; see, e.g., \citet{chenLargeSampleSieve2007} or \citet{picardDensityEstimationModel2007} for the full chaining/peeling details.

\paragraph{Conclusion.}
Hence (I) + (II) $= O_p(s_n \log n / n)$, and (III) $\le \inf_{f \in \cT_{s_n}}[R(f)-R(\eta_0)] + o(1)$. The extra $\lambda_n s_n$ is of the same order. Taking the constant $C$ in the lemma large enough, we obtain $R(\hat\eta) - R(\eta_0) \le C \bigl( \inf_{f \in \cT_{s_n}}[R(f)-R(\eta_0)] + \frac{s_n \log n}{n} \bigr) + o_p(1)$.

\paragraph{Remark: $\varepsilon_n$-optimal minimizers.}
We say $\hat\eta$ is \emph{$\varepsilon_n$-optimal} if its penalized empirical risk is at most $\inf_{f \in \cT_{M_n}}(\cdots) + \varepsilon_n$. Then $R_n(\hat\eta) - R_n(f_n) \le \lambda_n s_n + \varepsilon_n$; the same oracle inequality holds provided $\varepsilon_n = o(s_n \log n / n)$ (e.g., $\varepsilon_n = o(1/n)$), with $\varepsilon_n$ absorbed into the $o_p(1)$ term. GOSDT-style solvers with an appropriate tolerance produce such an $\hat\eta$.
\end{proof}

\begin{lemma}[Optimal choice of $s_n$]
\label{lem:optimal-sn}
With $s_n$ as the rate index (the penalized minimizer has $\hat{s} = O_p(s_n)$ leaves by Lemma~\ref{lem:hats}), balance the approximation term $s_n^{-2\beta/d}$ and the estimation term $s_n (\log n)/n$. Setting $s_n \asymp n^{d/(2\beta+d)}$ yields excess risk $R(\hat\eta) - R(\eta_0) = O_p(n^{-2\beta/(2\beta+d)})$. By the loss--norm link, $\|\hat\eta - \eta_0\|_{L^2(P)}^2 \le \phi(O_p(n^{-2\beta/(2\beta+d)})) = O_p(n^{-2\beta/(2\beta+d)})$ (e.g., if $\phi(\delta) \lesssim \delta$), so
\[
\|\hat\eta - \eta_0\|_{L^2(P)} = O_p\bigl( n^{-\beta/(2\beta+d)} \bigr).
\]
This is the classical minimax rate for $H^\beta([0,1]^d)$.
\end{lemma}
\begin{proof}[Proof idea]
Set $s_n^{-2\beta/d} \asymp s_n (\log n)/n$ and solve for $s_n$; substitute into the oracle inequality to get the excess-risk rate; apply the loss--norm link to obtain the $L^2(P)$ rate.
\end{proof}

\begin{theorem}[Nuisance rate]
\label{thm:nuisance-rate}
Suppose $\eta_0 \in H^\beta([0,1]^d)$, $X \in [0,1]^d$, the loss--norm link holds, and $\hat\eta$ is the penalized ERM~\eqref{eq:perm} (loss plus $\lambda_n \times \#\text{leaves}$) over $\cT_{M_n}$ with $\lambda_n \asymp (\log n)/n$ and $s_n \asymp n^{d/(2\beta+d)}$. Then $\|\hat\eta - \eta_0\|_{L^2(P)} = O_p(n^{-\beta/(2\beta+d)})$. If $\beta/(2\beta+d) > 1/4$, then this rate is $o_p(n^{-1/4})$.
\end{theorem}
\begin{proof}
Combine Lemmas~\ref{lem:approx}--\ref{lem:optimal-sn} and the loss--norm link.
\end{proof}

\begin{corollary}[DML validity]
\label{cor:dml}
If each nuisance estimator satisfies $\|\hat\eta - \eta_0\|_{L^2(P)} = o_p(n^{-1/4})$ in the norms that appear in the orthogonal score expansion, then the remainder in the DML expansion is $o_p(n^{-1/2})$. Under identification (overlap, bounded moments) and $K$-fold cross-fitting, $\sqrt{n}(\hat\theta - \theta_0) \leadsto N(0, \sigma^2)$, where $\sigma^2$ is the variance of the influence function. No new proof is required; we invoke \citet{chernozhukov2018} and verify their conditions (overlap, boundedness, cross-fitting).
\end{corollary}

\subsection{DML with alternative nuisances (general theorem)}
\label{app:alt-nuisances}

The following result is of independent interest for DML with cross-fitting: it states that any fold-specific nuisance estimators that are ``close'' to the true nuisances in the sense that controls the DML remainder still yield $\sqrt{n}$-consistent, asymptotically normal inference. No structure (e.g., trees) is assumed; the result is generic.

\begin{definition}[$o_p(n^{-1/4})$-close nuisances]
\label{def:close}
For each fold $k \in \{1,\ldots,K\}$, let $\tilde{\eta}^{(-k)} = (\tilde{e}^{(-k)}, \tilde{m}_0^{(-k)}, \tilde{m}_1^{(-k)})$ be a nuisance estimator trained on data not in fold $k$. We say the $\tilde{\eta}^{(-k)}$ are \emph{$o_p(n^{-1/4})$-close} to $\eta_0$ if, for each $k$, we have $\|\tilde{e}^{(-k)} - e_0\|_{L^2(P)} = o_p(n^{-1/4})$, $\|\tilde{m}_0^{(-k)} - m_{0,0}\|_{L^2(P)} = o_p(n^{-1/4})$, and $\|\tilde{m}_1^{(-k)} - m_{1,0}\|_{L^2(P)} = o_p(n^{-1/4})$. Under the overlap condition ($0 < e_0(X) < 1$ bounded away from 0 and 1), these componentwise $L^2(P)$ norms are the relevant norms for the ATT score expansion; equivalently, $\max_{k=1,\ldots,K} \|\tilde{\eta}^{(-k)} - \eta_0\|_{L^2(P)} = o_p(n^{-1/4})$ when $\|\tilde{\eta}^{(-k)} - \eta_0\|_{L^2(P)}$ denotes the maximum of the three component norms.
\end{definition}

\begin{theorem}[DML with alternative nuisances]
\label{thm:dml-alt}
Assume $K$ is fixed. Under the same identification, overlap, and moment conditions as in \citet{chernozhukov2018} (unconfoundedness, overlap, bounded outcomes and propensity), suppose that for each fold $k$, $\tilde{\eta}^{(-k)}$ is trained on data not in fold $k$ and is $o_p(n^{-1/4})$-close to $\eta_0$ in the sense of Definition~\ref{def:close}. Let $\hat\theta$ be the DML estimator defined by
\[
\frac{1}{n} \sum_{i=1}^n \psi\bigl(O_i; \hat\theta, \tilde{\eta}^{(-k(i))}\bigr) = 0.
\]
Then $\sqrt{n}(\hat\theta - \theta_0) \leadsto N(0, \sigma^2)$, where $\sigma^2$ is the variance of the influence function $\psi(O; \theta_0, \eta_0)$.
\end{theorem}
\begin{proof}
By \citet{chernozhukov2018}, the only nuisance-dependent condition for $\sqrt{n}$-consistency and asymptotic normality is that the second-order remainder in the linearization of the moment condition be $o_p(n^{-1/2})$. For the ATT orthogonal score \eqref{eq:score}, the remainder is bounded in absolute value by a constant times products of the form $\|\tilde{e}^{(-k)} - e_0\|_{L^2(P)} \|\tilde{m}_a^{(-k)} - m_{a,0}\|_{L^2(P)}$ (and similar cross-terms), possibly with overlap-weighted norms; see \citet{chernozhukov2018}, Section~3.2 and their Assumption 3.2. Under overlap, the weights are bounded, so componentwise $\|\tilde{\eta}^{(-k)} - \eta_0\|_{L^2(P)} = o_p(n^{-1/4})$ implies that each such product is $o_p(n^{-1/2})$. Hence the remainder is $o_p(n^{-1/2})$. The remaining conditions in \citet{chernozhukov2018} (identification, cross-fitting with $K$ fixed, bounded moments) do not depend on the choice of nuisance estimator. Thus their conclusion applies with $\tilde{\eta}^{(-k)}$ in place of $\hat\eta^{(-k)}$.
\end{proof}

\subsection{Rashomon set per fold and intersection}
\label{app:rashomon}

We give formal definitions of the Rashomon set for each cross-fitting fold and of the intersection across folds. These are used in Section~\ref{sec:rashomon}. For a fixed nuisance (e.g., propensity $e$ or outcome regression $m_0$), the training data for fold $k$ is $\mathcal{D}^{(-k)}$ (all observations not in fold $k$); $n_{-k} = |\mathcal{D}^{(-k)}|$. Let $R_n^{(-k)}(f) = (1/n_{-k}) \sum_{i \notin \text{fold } k} L(Y_i, f(X_i))$ denote the empirical risk on $\mathcal{D}^{(-k)}$ for the relevant loss $L$, and let $\text{best}^{(-k)} = \inf_{f \in \cT_{M_n}} [ R_n^{(-k)}(f) + \lambda_n \cdot (\#\text{leaves of } f) ]$ be the minimal penalized risk on that fold.

\begin{definition}[Rashomon set per fold]
\label{def:rashomon-fold}
For $\varepsilon \ge 0$, the \emph{Rashomon set} for fold $k$ is
\[
R_k(\varepsilon) = \Bigl\{ f \in \cT_{M_n} : R_n^{(-k)}(f) + \lambda_n \cdot (\#\text{leaves of } f) \le (1+\varepsilon) \cdot \text{best}^{(-k)} \Bigr\}.
\]
Optionally one restricts to trees with $\#\text{leaves} \le s_n$ so that the set is compatible with the sieve. We say trees in $R_k(\varepsilon)$ are \emph{$(1+\varepsilon)$-near-optimal} on fold $k$.
\end{definition}

\begin{lemma}[Oracle inequality for near-minimizers]
\label{lem:oracle-near}
Under the same conditions as Lemma~\ref{lem:oracle} (with $R_n^{(-k)}$, $\mathcal{D}^{(-k)}$, and $\text{best}^{(-k)}$ in place of $R_n$, $n$, and the fold-$k$ minimizer). For any $f \in R_k(\varepsilon_n)$ with $\#\text{leaves}(f) \le s_n$ (or $\#\text{leaves}(f) = O_p(s_n)$ under the same $\lambda_n$ logic as Lemma~\ref{lem:hats}),
\[
R(f) - R(\eta_0) \le C \Bigl( \varepsilon_n \cdot \text{best}^{(-k)} + \frac{s_n \log n}{n} + \inf_{f' \in \cT_{s_n}} \bigl[ R(f') - R(\eta_0) \bigr] \Bigr) + o_p(1).
\]
\end{lemma}
\begin{proof}[Proof idea]
Decompose $R(f) - R(\eta_0)$ as $(R(f) - R_n^{(-k)}(f)) + (R_n^{(-k)}(f) - R_n^{(-k)}(f^*)) + (R_n^{(-k)}(f^*) - R(f^*)) + (R(f^*) - R(\eta_0))$ for $f^* \in \cT_{s_n}$ with $R(f^*) - R(\eta_0)$ near $\inf_{f' \in \cT_{s_n}}[R(f') - R(\eta_0)]$. Since $f \in R_k(\varepsilon_n)$, we have $R_n^{(-k)}(f) + \lambda_n \cdot (\#\text{leaves of } f) \le (1+\varepsilon_n) \cdot \text{best}^{(-k)}$, so $R_n^{(-k)}(f) - R_n^{(-k)}(f^*) \le \varepsilon_n \cdot \text{best}^{(-k)} + \lambda_n s_n + O(s_n \log n / n)$ (using that the minimizer has penalized risk $\text{best}^{(-k)}$). The empirical-process terms are $O_p(s_n \log n / n)$ as in Lemma~\ref{lem:oracle}. The bound follows.
\end{proof}

\begin{definition}[Intersection of Rashomon sets (structure-based)]
\label{def:intersection}
The \emph{intersection} of the Rashomon sets across folds is the set of \emph{tree structures} (splits) such that for every fold $k$, when that structure is refit on $\mathcal{D}^{(-k)}$ (i.e., leaf values are optimized on data not in fold $k$), the resulting tree belongs to $R_k(\varepsilon)$. Here, \emph{refit} means: minimize $R_n^{(-k)}(f) + \lambda_n \cdot (\#\text{leaves of } f)$ over all trees $f$ with that same structure (same splits, variable leaf values). Equivalently, a structure is in the intersection if it appears in $R_k(\varepsilon)$ for each $k$ when each fold's trees are fitted on $\mathcal{D}^{(-k)}$.
\end{definition}

\begin{proposition}[Rashomon selection and DML validity]
\label{prop:rashomon-dml}
Suppose that for each fold $k$, we choose $\tilde{\eta}^{(-k)} \in R_k(\varepsilon_n)$ (e.g., any tree in the fold-$k$ Rashomon set, or the refit of an intersection structure on $\mathcal{D}^{(-k)}$). A \emph{sufficient condition} for $\|\tilde{\eta}^{(-k)} - \eta_0\|_{L^2(P)} = o_p(n^{-1/4})$ is: $\varepsilon_n \to 0$ and $\varepsilon_n = o(n^{-1/2})$ (so that $\varepsilon_n \cdot \text{best}^{(-k)} = o_p(n^{-1/2})$ since $\text{best}^{(-k)} = O_p(1)$), under the same sieve and rate conditions as Theorem~\ref{thm:nuisance-rate}. Then the DML estimator $\hat\theta$ defined with $\tilde{\eta}^{(-k)}$ satisfies $\sqrt{n}(\hat\theta - \theta_0) \leadsto N(0, \sigma^2)$ by Theorem~\ref{thm:dml-alt}.
\end{proposition}
\begin{proof}
By Lemma~\ref{lem:oracle-near}, any $f \in R_k(\varepsilon_n)$ with $\#\text{leaves}(f) \le s_n$ satisfies $R(f) - R(\eta_0) \le C( \varepsilon_n \cdot \text{best}^{(-k)} + s_n \log n / n + \inf_{f' \in \cT_{s_n}}[R(f') - R(\eta_0)] ) + o_p(1)$. Under the theorem's conditions, the approximation term and $s_n \log n / n$ are $o_p(n^{-1/2})$; with $\varepsilon_n = o(n^{-1/2})$, we obtain $R(f) - R(\eta_0) = o_p(n^{-1/2})$. By the loss--norm link, $R(f) - R(\eta_0) \le \delta$ implies $\|f - \eta_0\|_{L^2(P)}^2 \le \phi(\delta)$; for squared loss $\phi(\delta) = \delta$, so excess risk $o_p(n^{-1/2})$ gives $\|f - \eta_0\|_{L^2(P)} = o_p(n^{-1/4})$. Thus $\tilde{\eta}^{(-k)}$ is $o_p(n^{-1/4})$-close to $\eta_0$, and Theorem~\ref{thm:dml-alt} applies.
\end{proof}

\subsection{High-dimensional caveat}
If the tree class is allowed to split on all of many covariates, the effective complexity can grow and the preceding rates may not hold. To retain the same rate, one can assume (a) \emph{sparsity} of the regression function (few relevant covariates), (b) \emph{restricted depth} so that the effective dimension in the tree is bounded, or (c) \emph{margin} or regularity conditions. Such assumptions are standard in nonparametrics.

% Bibliography
\bibliographystyle{apalike}
\bibliography{dmltree-refs}

\end{document}
